% GENERATED FILE. Edit canonical lessons, then run npm run generate:books.
% canonical-source-hash: fnv1a64:41ad59a3df5c3543
% canonical-lessons: TE-C212-tirika, TE-C212-ventane, TE-C212-civaraku, TE-C212-varaku, TE-C212-sankranti

\chapter{Free time, Immediately, Finally, Until, Sankranti}
\label{ch:te-free-time-immediately-finally-until-sankranti}
\hlchaptermodality{212}

\begin{chapteropening}
\textbf{By the end of this chapter:} \emph{I can say free time, immediately, finally, until and Sankranti.}

Complete the last lesson of chapter 212: I can say free time, immediately, finally, until and Sankranti.
\end{chapteropening}

\section[\texorpdfstring{\te{తీరిక} (tīrika)}{తీరిక (tīrika)}]{\te{తీరిక} (tīrika) — free time, leisure}

\label{lesson:TE-C212-tirika}



\begin{quote}
Before the new one: say the Telugu for sunstroke, then the Telugu for steam.
\end{quote}

\subsection*{You'll want to know: \te{తీరిక}}
\textbf{\te{తీరిక}} — \emph{tīrika} — \textquotedblleft{}free time, leisure\textquotedblright{}.

\begin{grammarlens}[title={What you've built}]
One more word for this chapter.
\end{grammarlens}

\subsection*{Your turn}
\begin{itemize}
\raggedright
  \item \emph{Say it:} \emph{tīrika}
  \item \emph{Say it:} \emph{tīrika}, once more
  \item \emph{Say it:} say \emph{tīrika}
  \item \emph{Recall:} say the Telugu for sunstroke, then the Telugu for steam, then say \emph{tīrika} again
\end{itemize}

\begin{tcolorbox}[breakable,colback=teal!4,colframe=teal!35!black,title={Before you move on}]
What does \te{తీరిక} mean? (\textquotedblleft{}Free time, leisure\textquotedblright{}.) Say it once more.
\end{tcolorbox}
\section[\texorpdfstring{\te{వెంటనే} (veṇṭanē)}{వెంటనే (veṇṭanē)}]{\te{వెంటనే} (veṇṭanē) — immediately}

\label{lesson:TE-C212-ventane}



\begin{quote}
Before the new one: say the Telugu for steam, then the Telugu for free time.
\end{quote}

\subsection*{You'll want to know: \te{వెంటనే}}
\textbf{\te{వెంటనే}} — \emph{veṇṭanē} — \textquotedblleft{}immediately\textquotedblright{}.

\begin{grammarlens}[title={What you've built}]
One more word for this chapter.
\end{grammarlens}

\subsection*{Your turn}
\begin{itemize}
\raggedright
  \item \emph{Say it:} \emph{veṇṭanē}
  \item \emph{Say it:} \emph{veṇṭanē}, once more
  \item \emph{Say it:} say \emph{veṇṭanē}
  \item \emph{Recall:} say the Telugu for steam, then the Telugu for free time, then say \emph{veṇṭanē} again
\end{itemize}

\begin{tcolorbox}[breakable,colback=teal!4,colframe=teal!35!black,title={Before you move on}]
What does \te{వెంటనే} mean? (\textquotedblleft{}Immediately\textquotedblright{}.) Say it once more.
\end{tcolorbox}
\section[\texorpdfstring{\te{చివరకు} (civaraku)}{చివరకు (civaraku)}]{\te{చివరకు} (civaraku) — finally, at last}

\label{lesson:TE-C212-civaraku}



\begin{quote}
Before the new one: say the Telugu for free time, then the Telugu for immediately.
\end{quote}

\subsection*{You'll want to know: \te{చివరకు}}
\textbf{\te{చివరకు}} — \emph{civaraku} — \textquotedblleft{}finally, at last\textquotedblright{}.

\begin{grammarlens}[title={What you've built}]
One more word for this chapter.
\end{grammarlens}

\subsection*{Your turn}
\begin{itemize}
\raggedright
  \item \emph{Say it:} \emph{civaraku}
  \item \emph{Say it:} \emph{civaraku}, once more
  \item \emph{Say it:} say \emph{civaraku}
  \item \emph{Recall:} say the Telugu for free time, then the Telugu for immediately, then say \emph{civaraku} again
\end{itemize}

\begin{tcolorbox}[breakable,colback=teal!4,colframe=teal!35!black,title={Before you move on}]
What does \te{చివరకు} mean? (\textquotedblleft{}Finally, at last\textquotedblright{}.) Say it once more.
\end{tcolorbox}
\section[\texorpdfstring{\te{వరకు} (varaku)}{వరకు (varaku)}]{\te{వరకు} (varaku) — until, up to}

\label{lesson:TE-C212-varaku}



\begin{quote}
Before the new one: say the Telugu for immediately, then the Telugu for finally.
\end{quote}

\subsection*{You'll want to know: \te{వరకు}}
\textbf{\te{వరకు}} — \emph{varaku} — \textquotedblleft{}until, up to\textquotedblright{}.

\begin{grammarlens}[title={What you've built}]
One more word for this chapter.
\end{grammarlens}

\subsection*{Your turn}
\begin{itemize}
\raggedright
  \item \emph{Say it:} \emph{varaku}
  \item \emph{Say it:} \emph{varaku}, once more
  \item \emph{Say it:} say \emph{varaku}
  \item \emph{Recall:} say the Telugu for immediately, then the Telugu for finally, then say \emph{varaku} again
\end{itemize}

\begin{tcolorbox}[breakable,colback=teal!4,colframe=teal!35!black,title={Before you move on}]
What does \te{వరకు} mean? (\textquotedblleft{}Until, up to\textquotedblright{}.) Say it once more.
\end{tcolorbox}
\section[\texorpdfstring{\te{సంక్రాంతి} (saṅkrānti)}{సంక్రాంతి (saṅkrānti)}]{\te{సంక్రాంతి} (saṅkrānti) — Sankranti, the January harvest festival}

\label{lesson:TE-C212-sankranti}



\begin{quote}
Before the new one: say the Telugu for finally, then the Telugu for until.
\end{quote}

\subsection*{You'll want to know: \te{సంక్రాంతి}}
\textbf{\te{సంక్రాంతి}} — \emph{saṅkrānti} — \textquotedblleft{}Sankranti, the January harvest festival\textquotedblright{}.

\begin{grammarlens}[title={What you've built}]
One more word for this chapter.
\end{grammarlens}

\subsection*{Your turn}
\begin{itemize}
\raggedright
  \item \emph{Say it:} \emph{saṅkrānti}
  \item \emph{Say it:} \emph{saṅkrānti}, once more
  \item \emph{Say it:} say \emph{saṅkrānti}
  \item \emph{Recall:} say the Telugu for finally, then the Telugu for until, then say \emph{saṅkrānti} again
\end{itemize}

\begin{tcolorbox}[breakable,colback=teal!4,colframe=teal!35!black,title={Before you move on}]
What does \te{సంక్రాంతి} mean? (\textquotedblleft{}Sankranti, the January harvest festival\textquotedblright{}.) Say it once more.
\end{tcolorbox}
